% !TEX root = ../main.tex
\chapter{Linearization and Numerical Specialization}

\section{Operating Point}

The control design operating point is the horizontal beam equilibrium with the ball at
rest at the center of the usable sensing window:

\begin{equation}
  r = r_0 = \SI{0.1256}{m}, \qquad
  \theta = 0, \qquad
  \dot{r} = 0, \qquad
  \dot{\theta} = 0
\end{equation}

At this operating point the actuator must balance gravity with

\begin{equation}
  \tau_0 = (mr_0 + M_bL_{\text{com}})g
\end{equation}

In the implemented cascaded controller, this equilibrium torque is handled implicitly
through the beam-angle servo and balance-zero procedure.

\section{Perturbation Variables}

Define small deviations from equilibrium by

\begin{align}
  r &= r_0 + x \\
  \theta &= \delta \\
  \dot{r} &= \dot{x} \\
  \dot{\theta} &= \dot{\delta}
\end{align}

Substituting these expressions into the nonlinear equations and neglecting second-order
products gives the linearized model.

\section{Linearized Ball Dynamics}

Starting from \eqref{eq:eom_r}:

\begin{equation}
  \frac{5}{3}m\ddot{x} + b_x\dot{x} + mg\delta = 0
\end{equation}

Dividing through by the effective rolling mass $m_e = \tfrac{5}{3}m$ gives

\begin{equation}
  \ddot{x} + \beta \dot{x} = -\alpha g \delta
  \label{eq:linear_ball}
\end{equation}

with

\begin{align}
  \alpha &= \frac{m}{m_e} = \frac{3}{5} = 0.6 \\
  \beta &= \frac{b_x}{m_e} = \frac{0.0125}{0.0046667} = \SI{2.6786}{s^{-1}}
\end{align}

Numerically, the ball-position plant becomes

\begin{equation}
  \ddot{x} + 2.6786\dot{x} = -5.886\,\delta
\end{equation}

or in transfer-function form,

\begin{equation}
  \frac{X(s)}{\Theta(s)} =
  \frac{-5.886}{s^2 + 2.6786s}
  = \frac{-5.886}{s(s + 2.6786)}
\end{equation}

\section{Linearized Beam Dynamics}

Linearizing \eqref{eq:eom_theta} gives

\begin{equation}
  (J + mr_0^2)\ddot{\delta}
  + b_{\theta}\dot{\delta}
  + (mr_0 + M_bL_{\text{com}})g\delta
  = \tau - \tau_0
  \label{eq:linear_beam}
\end{equation}

Substituting the evaluated parameters yields

\begin{align}
  J_{\text{total}}(r_0)
    &= J + mr_0^2
     = \SI{8.157e-4}{kg.m^2} \\
  (mr_0 + M_bL_{\text{com}})g
    &= \SI{0.04804}{N.m/rad}
\end{align}

Ignoring damping for an estimate, the beam mode has natural frequency

\begin{equation}
  \omega_{\text{beam}}
  = \sqrt{\frac{0.04804}{8.157\times 10^{-4}}}
  = \SI{7.67}{rad/s}
  \approx \SI{1.22}{Hz}
\end{equation}

\section{Specialized Parameters}

Table~\ref{tab:specialized_parameters} collects the key values used for control
interpretation.

\begin{table}[tbp]
  \centering
  \small
  \caption{Numerically specialized parameters around the nominal operating point.}
  \label{tab:specialized_parameters}
  \begin{tabularx}{\linewidth}{@{}l c X@{}}
    \toprule
    Symbol / quantity & Value & Meaning \\
    \midrule
    $m_e$ & \SI{0.004667}{kg} & Effective rolling mass of the hollow ball \\
    $\alpha$ & 0.6 & Ball-to-effective-mass ratio \\
    $\beta$ & \SI{2.6786}{s^{-1}} & Linearized viscous coefficient in the ball equation \\
    $\alpha g$ & \SI{5.886}{m/s^2} & Beam-angle to ball-acceleration gain \\
    $J_{\text{total}}(r_0)$ & \SI{8.157e-4}{kg.m^2} & Total beam-side rotational inertia at center position \\
    $(mr_0 + M_bL_{\text{com}})g$ & \SI{0.04804}{N.m/rad} & Linearized gravitational restoring torque coefficient \\
    $\omega_{\text{beam}}$ & \SI{7.67}{rad/s} & Approximate beam natural frequency without extra actuator dynamics \\
    \bottomrule
  \end{tabularx}
\end{table}

\section{Implication for Controller Structure}

The linearized ball equation has an integrator and a stable real pole, while the beam
subsystem is comparatively faster and directly actuated. That separation motivates the
implemented cascaded structure:

\begin{itemize}
  \item an inner loop regulates beam angle using the AS5600 measurement and stepper
  actuation, and
  \item an outer loop commands beam angle from ball-position error, velocity, and a
  constrained integral term.
\end{itemize}

The actual firmware is more detailed than the ideal linear model. It retains
quantization, deadband, validity gating, recovery logic, and empirical limits, which
are described in the next chapter.
